\subsection{Observations and Design Implications}
\label{sec:design-insights}

\paragraph{Observation 1: recovery speed has workflow-dependent value.}
The fan-out example in Section~\ref{sec:motivation-workflow} shows that
the same recovery delay can be absorbed on one branch and delay the join
on another. A branch-duration gap is not itself the slack at failure time.
This suggests comparing the delay a request can absorb with the recovery
cost of each storage tier before allocating capacity. AegisKV estimates
slack from the workflow DAG and uses it for protection tiers and
competition order (Section~\ref{sec:workflow-protection}).

\paragraph{Observation 2: access is concentrated within a request.}
We collect actual selector indices from MiniCPM4-8B, NOSA-1B, and
NOSA-3B on three PG19 documents per model. Each capture has a 32K-token
history and a causal continuation; 256- and 512-query windows are nested
within the same capture. Figure~\ref{fig:design-concentration} ranks
\emph{(layer, 64-token block)} objects, each containing all KV heads,
by selection frequency. The hottest 20\% of historical KV bytes account
for 92.2\%/90.3\% of selection events for NOSA-1B,
91.5\%/89.5\% for NOSA-3B, and 69.6\%/66.0\% for MiniCPM4-8B
in the two windows. The two NOSA curves are close in these configurations;
this does not establish independence from model scale.

\begin{figure*}[t]
  \centering
  \includegraphics[width=\textwidth]{figures/design-sparse-concentration.pdf}
  \caption{Concentration of native sparse-selection events in historical KV.
    Each model has three PG19 documents and one causal capture per document;
    the 256/512-query windows are nested, not independent runs.
    Curves show the pointwise median across documents after offline
    frequency ranking. Repeated selections of a block by different queries
    or selection groups count repeatedly. The denominator excludes new KV.
    Event coverage is neither an online hit rate nor a recovery-byte saving.}
  \label{fig:design-concentration}
  \Description{The highest-frequency layer-block objects cover more selector
    events than uniform selection; the degree of concentration differs by model.}
\end{figure*}

This concentration motivates selective protection, but does not determine
how much KV recovery requires: a block selected repeatedly is transferred
only once. AegisKV therefore uses historical hotness to predict a needed
set, while retaining a complete path for prediction misses
(Section~\ref{sec:kv-placement}).

\paragraph{Observation 3: Peer HBM removes exposed transfer wait.}
Once a committed DRAM copy exists, recovery can already avoid recomputing
the entire history. The additional value of Peer HBM is the transfer time
that remaining computation cannot hide. During layer-by-layer recovery,
earlier layers need their KV sooner, while later preloads can overlap
earlier computation. This motivates protecting a front-layer prefix of
hot sets until preload wait is eliminated. The calibrated pipeline in
Section~\ref{sec:kv-placement} makes this argument explicit; it does not
establish optimality over arbitrary placements.

These observations leave two systems challenges: converting layer/block
choices into a bounded, replaceable placement plan, and executing that
plan when the recovery target still serves foreground requests.
AegisKV uses whole-footprint admission to bound placement
(Section~\ref{sec:capacity-allocation}). Independent commits and
per-request readiness govern recovery
(Sections~\ref{sec:parallel-recovery} and~\ref{sec:kv-runtime}).
